\section{Discussion}
% 1. summarise what we have shown.
We created a sex-balanced virtual cohort of 50 whole-heart models derived from clinical CT scans of 
both healthy individuals and \HF\ patients. 
We used a semi-automated pipeline to generate these models,
which includes segmentation, mesh generation, fibre orientation assignment, and simulation setup.
The models enabled the exploration of differences in volume, and activation time across sex and disease groups, 
using benchmark electrophysiology and mechanics simulations. 
With this study, we demonstrate that anatomical differences between sexes and across disease states have a 
significant impact on both electrical and mechanical simulation predictions in the heart. 
This highlights the critical need to account for such variations when developing and 
applying computational models in cardiology.

% 2. comment on other workflows
A variety of computational pipelines have been proposed to construct heart models, 
ranging from atria-only or ventricles-only representations to anatomically detailed four-chamber reconstructions.
The works by~\cite{roney-2023-atriamodels} and~\cite{azzolin-2023-atriamodels} proposed 
patient-specific atrial models, including anatomical regions and fibre orientations from clinical data.
In contrast, other pipelines focus only on ventricular chambers~\cite{govil-2023-bivmodels,banerjee-2021-bivmodels}.
A notable whole-heart model was developed by~\cite{zhang-2012-fourchmodels}, 
which builds detailed four-chamber meshes of the entire heart from patient imaging and has been enhanced 
since to improve 
simulation fidelity, 
multi-scale integration, and 
clinical translation.
The resulting models are anatomically detailed, but the method depended on extensive user guidance to 
create the initial atlas, limiting throughput and scale.
\replaced[id=JASL]{
    Other groups, such as Gillette et al.~\cite{gillette_personalized_2022}, have also 
    demonstrated personalised whole-heart electrophysiology models, but again with 
    significant manual setup and typically small cohorts.
}{
    Other groups~\cite{gillette_personalized_2022} have also demonstrated personalised 
    whole-heart electrophysiology models, but again with significant manual setup and 
    typically small cohorts.
}
In summary, existing pipelines either produce atria-only or ventricle-only models, 
or else deliver high-detail four-chamber meshes at the cost of manual effort.
Our framework seeks to overcome these limitations by producing anatomically complete four-chamber 
(atria + ventricles) models with minimal manual input, combining scalability with high anatomical fidelity.

\subsection{Building Patient-Specific Whole-Heart Models: Workflow Insights}
The models were created using a standardised pipeline, CEMRG heartbuilder, 
which allows for the generation of patient-specific models from CT scans. 
The segmentation framework has enabled a substantial reduction in overall model generation time relative to 
earlier methods. The pipeline requires approximately 5 hours per heart on average.
It is also significantly more detailed than most publicly available pipelines. 
For instance, methods evaluated in the MICCAI 2024 Whole Heart Segmentation Challenge~\cite{miccai2024_challenge} 
typically produce 7–10 output labels, insufficient for physiologically realistic biophysical simulations. 
In contrast, our approach employs a 37-label standard~\cite{strocchi2020_cohort,strocchi2020_simulating}, 
encompassing detailed delineation of blood pools, myocardia, valve planes,
and major vessels across all four chambers.
\added[id=JASL]{
    The current implementation applies uniform meshing parameters globally 
    across the whole heart. While regional refinement could theoretically be 
    achieved by meshing structures separately and stitching at interfaces, 
    we prioritized mesh integrity to ensure numerical stability in coupled 
    electromechanical simulations, where interface discontinuities can 
    introduce artifacts.
}
\added[id=JASL]{
    The strong correlations between geometric metrics and simulation outputs 
    (Supplementary Table~S2, Figure~S2) demonstrate that anatomical variability 
    is a primary driver of functional differences when using standardized 
    electrophysiological and mechanical parameters. Chamber volume alone explained 
    92\% of variance in volume change during inflation ($\rho = 0.96$), indicating 
    that anatomical size dominates mechanical behavior under uniform constitutive 
    laws. Similarly, mesh element count (a proxy for anatomical size) predicted 
    87\% of variance in right ventricular activation time, consistent with 
    geometric path length determining conduction duration.

    These findings validate our approach of isolating anatomical effects through 
    parameter standardization. However, they also highlight the limitation: 
    real-world patient-specific predictions would require calibration of tissue 
    properties (contractility, stiffness, conduction velocity) to observed 
    clinical metrics.
}

% \subsection{Discussion on the results}
\subsection{Impact of Anatomical Variability on Electromechanical Predictions}
All whole-heart models were simulated using identical parameters for both electrophysiology and mechanics. 
% This design choice ensured that any inter-group variation in simulation output arose solely from 
% differences in anatomy, rather than from parameter fitting or model calibration. 
We did not include active contraction in the benchmark as these simulations often require patient specific 
parameters to match contraction, pre-load and after load, 
reducing our ability to compare consistent simulations across anatomies.
As such, the results reflect the isolated impact of structural variation due to sex or disease status on 
predicted electrical activation and mechanical inflation.
The results from the two-way ANOVA showed, in most cases, non significant interactions.
This means that we could only look at the independent effects: control vs \HF\ and 
male vs female. 
The characterisation of effect sizes (small, moderate, and large) by~\cite{gignac-2016-effects}, 
were included as a guideline, based on the distribution of effect sizes obtained.

\subsection{Comparison with Clinical Reference Ranges}
Validation against population-based MRI reference data~\cite{petersen-2017-reference} 
confirmed that control models reproduced expected sex-specific chamber sizes.
Male and female \LV\ and \RV\ volumes were largely consistent with normal ranges, 
supporting the accuracy of our anatomical pipeline.
In heart failure, \LV\ enlargement was most pronounced in males, frequently exceeding 
reference thresholds, consistent with classical patterns of systolic dilation.
Females also showed increased \LV\ volumes, though these remained mostly within upper limits, 
highlighting sex differences in disease expression.
\RV\ enlargement was less prominent in both males and females with \HF. 
This might be reflecting clinical reports that \RV\ size is less sensitive to early disease~\cite{ro-2023-rv_volume}.
These results demonstrate that the virtual cohort captures both healthy anatomy and pathological remodelling, 
while also exposing the limitations of fixed reference ranges that do not account for sex and disease interplay.

\paragraph{Impact of Sex and Heart Failure in \emph{In-Silico} Studies}
Our simulations consistently separated groups by sex and disease status, even in the absence of tuned 
electrophysiological or mechanical parameters.
Absolute \LV\ volume increases were particularly marked in male \HF, reproducing the clinical signature of 
pathological dilation~\cite{kasa-2024-lv_volume,ito-2021-lv_volume}.
Larger $\deltaVol{LV}$ in \HF\ patients suggest structurally thinner or more distensible 
walls~\cite{zile-2004-stiffness,aurigemma-2006-stiffness}.
Similarly, larger $\deltaVol{LA}$ would indicate other issues like increased prediction of mortality or  
hospitalisation~\cite{Carpenito-2021-la-stiffness,kim-2023-la-stiffness}.
When expressed as a normalised change, however, \HF\ females showed the greatest relative differences, 
a finding likely driven by smaller baseline chamber sizes.
This emphasises that absolute and relative measures provide complementary insights and must be interpreted carefully 
in sex-stratified analyses~\cite{steuernagel2023_sex_bias}.
Electrical activation followed similar logic: ventricular and atrial activation were prolonged in \HF\ and in males, 
consistent with clinical QRS and conduction observations~\cite{sanders-2003-atrialremodelling,buck-2008-crt}, 
though interactions were not significant.
These convergent trends across independent metrics highlight that anatomical variability alone can generate 
clinically relevant patterns of electromechanical behaviour, a finding consistent with model-based predictions 
of sex-specific conduction behaviour~\cite{Lee2019_gender}.
These findings suggest that even without personalisation of tissue properties or conduction parameters, 
anatomically driven whole-heart simulations can outline group-level trends in cardiac function.
The consistent importance of sex and disease differences across multiple metrics reinforces the clinical 
relevance of the modelling pipeline and motivates its use in larger-scale studies~\cite{strocchi2020_cohort}.

\subsection{Establishing Model Credibility and Validation Strategies for In Silico Clinical Trials}
In silico studies using virtual cohorts of patients, referred to as in silico clinical trials (ISCTs), 
require careful demonstration of model credibility for their conclusions to be considered reliable. 
Credibility is primarily established by performing validation, that is, comparison of model predictions 
with real-world data. 
However, unlike relatively simple models of standalone medical devices, or drug safety/efficacy studies 
using a generic action potential model, ISCTs present a much wider range of possible validation strategies. 
As discussed in detail in~\cite{pathmanathan-credibility-2024}, possible evaluation activities include: 
validation of a device/drug model in isolation 
(e.g., bench test validation of a pacemaker lead mechanics model); 
individual-level validation of patient models in the cohort 
(i.e., comparing predictions from a specific patient model with clinical data from the specific patient 
the model represents), 
population-level validation of the cohort or sub-cohorts 
(e.g., comparing the distribution of model outputs across the cohort with clinical data from the 
corresponding real-world population), 
validation of coupled device-patient or drug-patient models 
(e.g., PKPD validation studies), 
and validation of any statistical model used to convert patient model outputs into clinical endpoints. 
Also, there are important additional activities that should be performed, such as assessing the virtual 
cohort representativeness; 
these are not strictly model validation activities and more akin to model input validation for 
simple models.
 
Overall, an ISCT is expected to be supported by a body of evidence on model credibility. 
For a virtual cohort such as the sex-balanced 50 patient cohort provided here, 
intended for re-use in variety of future device or drug ISCTs, 
a hierarchical evaluation approach can be used to generate an initial body of credibility evidence 
to support a future ISCT.
In this work, we have performed the first steps of such a hierarchical approach, 
by assessing cohort representativeness (and simultaneously assessing accuracy of the model generation pipeline) 
through comparison the ventricular and atrial volumes in the male and female cohort subsets with 
corresponding UK BioBank data.
 
A natural next step would be to perform sex-specific population-level validation: 
demonstrating that the distribution of model outputs, for each of the male and female sub-cohorts, 
match the distribution observed in the male and female populations, 
particularly for outputs for which a statistically significant sex difference was observed. 
For instance, predicted chamber $\TAT{}$ distributions could be compared with sex-specific population 
QRS duration data, 
after an appropriate transformation to account for the fact that $\TAT{}$ is only a surrogate for QRS.
Individualized validation, while more challenging, could be performed for newly generated patient models.
Ultimately, these results would be supplemented by context-specific validation when the cohort is used in a 
device/drug ISCT, together with other credibility assessment activities such as verification and 
uncertainty quantification~\cite{health_assessing_2023}.

\subsection{Limitations}
This work has four main limitations that should be considered when interpreting the results.
% + sample size 
First, the small sample size within each category of \HF , narrow and wide QRS duration, 
limits the statistical power for detecting subtle effects or interactions. 
Nonetheless, by balancing sex and \HF\ subtype, our dataset offers a structured 
starting point for exploring sex-disease interactions in-silico.
Furthermore, it can serve as a foundation for synthetic cohort 
generation~\cite{rodero-2021-virtualcohort1000,rodero-2021-virtualcohortextreme}
with a more diverse range of anatomical and clinical characteristics. 
% + complex simulations
Secondly, the simulations performed in this work were benchmark simulations 
% + parameters not calibrated 
and the parameters were not calibrated to the individual anatomies.
However, the use of a common set of parameters across all cases
ensured that the differences observed in the results were purely anatomical.
% time spent on the pipeline
Third, the pipeline requires 5 hours per heart on average, 
with 2 hours spent performing manual corrections required at the segmentation stage.
Even so, this is a significant reduction compared to previous methods.
Finally, limitations remain at the segmentation stage. 
Because this step underpins all subsequent model construction, 
we discuss its implications in more detail in the following section.

\added[id=JASL]{
    Our simulations employed uniform electrophysiological and 
    mechanical parameters across all models to isolate the effects of anatomical 
    variability. While this approach successfully captured anatomical influences 
    on cardiac function, evidenced by the strong correlation between simulated 
    activation times and QRS duration in control hearts, $r=0.649, p<0.005$, it 
    does not account for patient-specific tissue properties such as fibrosis or 
    conduction system disease. The absence of correlation in heart failure patients 
    ($r=-0.127, p=0.64$) reflects this limitation and highlights the need for 
    patient-specific parameter calibration in future work.
}

\subsubsection{Critical Role of Segmentation in Model Accuracy}
The segmentation stage is a critical bottleneck in the pipeline,
requiring notable manual input to ensure anatomical fidelity.
The most frequent issues involve small gaps or discontinuities in the myocardium, 
where the wall becomes excessively thin.
The segmentation method proposed by~\cite{hao2021_ccta} is not robust against some artefacts
introduced by implanted devices, such as the bright streaks caused by leads or stents in \HF\ patients. 
Manual steps and corrections in the pipeline could potentially be replaced with more 
robust algorithms. 
For example, recent work by~\cite{qayyum-2025-foundationmodel} demonstrated a 
foundation model capable of deriving full myocardial masks without user input.
Finally, given the limited soft-tissue contrast and resolution of standard CT scans, 
the myocardial thickness in chambers outside of the \LV\ is set to a user-specified constant.
This approximation may not reflect the true myocardial thickness 
and may underestimate simulation variability.

Despite these limitations, our current approach balances physiological realism with scalability, 
enabling whole-heart simulations across a sizeable cohort. 
Incorporating more refined fibre estimation methods remains an important direction for improving 
predictive accuracy in personalised cardiac models.
By incorporating this level of anatomical detail into a reproducible, Python-based workflow, 
our pipeline bridges the gap between clinical image segmentation and simulation-ready models. 
It provides a viable route to building whole-heart virtual cohorts for in-silico trials and 
aligns with regulatory guidelines on model credibility and anatomical completeness~\cite{fda2023_web}.

